\documentclass[twocolumn,9pt]{article}
%\documentclass[9pt]{article}
%\documentclass[../main.tex]{subfiles}
\usepackage{amsmath}
\usepackage{amssymb}
\usepackage{graphicx}
\usepackage{amsthm}
\usepackage{breqn}
\usepackage{float}
\usepackage{enumerate}
\usepackage{comment}
\usepackage{xcolor}
\usepackage{hyperref}

\newcommand\norm[1]{\lVert#1\rVert}
\newtheorem{theorem}{Theorem}[section]
\newtheorem{corollary}{Corollary}[theorem]
\newtheorem{proposition}{Proposition}[theorem]
\newtheorem{lemma}[theorem]{Lemma}
\theoremstyle{definition}
\newtheorem{definition}{Definition}[section]
\theoremstyle{remark}
\newtheorem*{remark}{Remark}

%\begin{figure}
%\includegraphics[width=1.0\textwidth]{plt_introduction.pdf}
%\caption{Two paths having small Euclidean distance, but big Signature
%            difference}
%\includegraphics[scale=0.5]{plt_introduction.pdf}
%\centering
%\end{figure}

%\title{The Signature Transform for Time Series Classification (beta3)}
\author {-}

\begin{document}

%\maketitle

\begin{abstract}
Use the signature to linearize the problem of time series prediction.
Use the bayesian theory to solve the problem taking advantage of
the linear formulation.
\end{abstract}

\begin{section}{The Linear Case}
Suppose to have a dataset of points $\{X_i, y_i\}$ connected by a linear
model perturbed by noise:
\begin{equation}
y_i = <L, X_i> + \eta
\end{equation}
for an appropriate array $L$, and $\eta$ Gaussian noise.
The goal is to find a good value for $L$, such that it is able to reproduce
all the data $y$ starting from $X$, for points in a training set,
as well as to predict correctly on unseen data from a target 
validation set.

We follow the classic Bayesian approach, which puts probability distributions
on $L$ insted of fixed, determined values.
Start with a Gaussian $N(0, I)$ on $L$: observing the data modifies the
distribution and allows to perform predictions in a probabilistic way.


Since the model $L$ is linear, there are close formulaes for that!
And, crucial: predictions will also follow a (dimensional)
Gaussian distributions, meaning that every time we have either
a mean (also more, i.e. most likely value) as well as 
confidence intervals here computed as $1.96 \sigma$.

\end{section}


\begin{section}{Preprocessing Data}
Giv

\end{section}

\begin{section}{The Linear Case}

\end{section}



%\bibliographystyle{alpha} % alternative is "plain" reference style
%\bibliography{./references/thebibliography}
\end{document}
